% Chapter 6. Every number is a macro (generated/numbers.tex) or a \needsclaim box.
\chapter{The Window Allocator}
\label{ch:allocator}

Chapters~\ref{ch:cost} and~\ref{ch:value} found that on this device the
silence is mostly the recognizer's, that the thinker's speculation costs
without paying, and that the listener's own speculation has a penalty that can
be stopped at its source. This chapter builds the rule that follows. The
listening window is allocated per turn, to whichever consumer can finish
before the endpoint and whose output survives it. The chapter then measures
what that rule achieves, and asks whether a planner could do better.

\section{Design}
\label{sec:design}

\textbf{Committing while listening.} A hypothesis worker, a separate process
running the \texttt{tiny} engine, decodes only the audio after the committed
point. A LocalAgreement-2 committer \cite{machacek2023whisperstreaming}
commits the prefix that two consecutive hypotheses agree on, at segment
granularity, and never commits a word inside the last 0.3\,s of the buffer.
At the endpoint the final engine decodes only the uncommitted tail.

\textbf{Pacing.} Hypotheses are self-paced. After a decode of length $d$, the
worker rests $(1/\delta - 1)\,d$ before the next, with duty $\delta = 0.6$ and
$d$ the measured length of the last decode. Cadence is an output of the rule,
not an input. A duty near 1 split no more turns than this; the bound is kept
because it gives the same outcome with fewer decodes.

\textbf{The feasibility gate.} A hypothesis is issued only if the decodes
still needed to commit its words, with their rests, fit in the expected
remaining speech. The expected remaining speech is a prior utterance length
minus the audio heard so far. The prior is a point value, derived at run start
as the median duration of the set's files. On short utterances the gate
abstains outright.

\textbf{Freezing the orphan.} A hypothesis still decoding at the endpoint is
frozen with \texttt{SIGSTOP}. It is resumed only when the recognizer is idle
after the reply has played, and its stale result is discarded by sequence
number. This removes the listener's overlap penalty without killing the
worker. It costs the energy of the resumed decodes, and a turn that opens
while the worker is frozen gets no hypotheses.

\textbf{The allocator.} For the thinker, at every partial transcript the
allocator:
\begin{enumerate}
\item estimates the fraction $\hat f$ heard, from the prior;
\item reads usability $p_{\mathrm{usable}}(\hat f)$ from the curve measured
  offline (Chapter~\ref{ch:value});
\item admits a draft of budget $B$ only if $p_{\mathrm{usable}}(\hat f) \ge
  p^\ast$ and $B$ tokens fit in the window at the measured decode rate.
\end{enumerate}
The listener's gate runs independently of it. Every decision is logged with
its inputs and its reason.

\section{The listener without the gate}
\label{sec:pass2}

The first test of the listener ran it on every turn, with a thread worker,
no gate and no freeze. Against the reactive pipeline without partial decodes,
it changed TTFA on multi-step problems by \PtwoCI{} ($n=\PtwoN$ items, as a
reduction): slower, with an interval that includes zero. Its final decode ran
above the solo model of its own tail
(\needsclaim{final over solo model, multi-step}{2026-10-06 --- PASS 2 SCORED}).
On the short development set it was clearly slower
(\needsclaim{dev-set TTFA regression}{same entry}), a property of the design
at short durations. The registered prediction failed on this pass. Two
mechanisms were proposed for the failure, the orphaned decode at the endpoint
and speculation on utterances too short to pay; the freeze and the gate are
the responses to them, and only the first mechanism was tested directly.

\section{Preemption}
\label{sec:preemption}

Three designs for stopping the orphaned hypothesis were tried, each
pre-registered before its run.
\begin{itemize}
\item \textbf{Kill and respawn.} Falsified as run
  (\needsclaim{final over solo model under KILL}{2026-10-06 --- P9 FALSIFIED
  AS RUN}). The design respawned the worker at every endpoint, and by
  elimination the respawn caused the slowdown.
\item \textbf{Kill and promote a hot spare.} Not testable. Its validity check,
  which voided a run in which a spawn overlapped any endpoint-to-final window,
  failed in every attempt. The first attempt's favourable result was withdrawn
  for that reason.
\item \textbf{Freeze.} Against letting the orphan race the final, freezing it
  cuts paired TTFA by \PnineTTFA{} ($n=\PnineN$ items). The final runs at
  \PnineRatio{} of the solo model of its tail. Three qualifications hold every
  time this figure is cited:
  \begin{itemize}
  \item the interval's lower end lies below the registered 150\,ms bar;
  \item the result holds under an amendment made before the run, and under
    the original validity rule it would not be testable;
  \item it depends on a 3\,s quiet period after each reply.
  \end{itemize}
\end{itemize}

\section{The gated, freezing listener}
\label{sec:pass2k}

\begin{figure}[t]
  \centering
  \includegraphics[width=\textwidth]{paper_fig5_pass2k.pdf}
  \caption{The gated, freezing listener (pass 2k) and the ungated, unfrozen
  one (pass 2), each against the reactive pipeline without partial decodes:
  TTFA reduction, word-error-rate difference and raw energy per turn.}
  \label{fig:th-listener}
\end{figure}

With the gate and the freeze, against the reactive pipeline without partial
decodes, the listener cuts the median paired TTFA on multi-step problems by
\PthreeRed{} ($n=\PthreeN$ items; Figure~\ref{fig:th-listener}). This meets
the prediction we registered for it, though the interval's lower end lies
below the registered 300\,ms bar. On the development set the gate issued no
hypothesis, at no detectable TTFA cost (\DevTTFA). On multi-step problems:
\begin{itemize}
\item \textbf{Commits.} The listener commits a median \WordsCommitted{} of
  the words before the endpoint ($n=\WordsN$ items). That is short of the
  0.60 we expected.
\item \textbf{Accuracy.} The word-error-rate cost is bounded but not shown to
  be absent: mean \PfourMean{} \PfourCI, with \PfourWorse{} items worse and
  \PfourBetter{} better.
\item \textbf{The final decode} runs at \PsixRatio{} of the solo model of its
  tail. Pass 2's final ran above its solo model, but the two passes differ
  in preemption, gate, quiet period and build, so this is not a paired
  comparison.
\item \textbf{Energy.} The listener costs \EnergyListenerr{} per turn, net of
  the idle floor, which is \SevenJms{} of time saved. Frozen decodes resumed
  and discarded after the reply cost a net mean of \StaleMean\,J each, and
  \StaleTotal\,J over all \StaleN.
\end{itemize}

\section{What the allocator decided}
\label{sec:decisions}

% Generated by src/scripts/paper_results.py. Do not edit.
\begin{table*}[t]
\centering\footnotesize
\caption{The allocator's decisions in one multi-step run, by reason: median (min to max) of the inputs each decision logged.}
\label{tab:allocator}
\begin{tabular}{lrrrrrrrr}
\toprule
reason & decisions & turns & f\_hat & p\_usable\_hat & feasible budget (tokens) & remaining (ms) & budget spent & drafts issued \\
\midrule
below\_break\_even & 537 & 81 & 0.391 (0.083 to 0.839) & 0.020 (0.013 to 0.137) & 96 (16 to 96) & 7878 (873 to 12690) & 0 & 0 \\
infeasible & 39 & 39 & 0.870 (0.846 to 0.935) & 0.163 (0.142 to 0.188) & 0 (0 to 0) & 385 (-630 to 771) & 0 & 0 \\
\bottomrule
\end{tabular}
\end{table*}


In one logged multi-step run the allocator took \AllocDecisions{} thinker
decisions (Table~\ref{tab:allocator}). \AllocBelow{} were below break-even:
the draft would have fit, but its usability at the estimated fraction was
below $p^\ast$. \AllocInfeasible{} were infeasible: usability had passed
break-even, but under the registered window model no draft fit. It spent the
window on the drafter \AllocSpends{} times. Every record carries the inputs
its reason rests on. The rule's constant output is therefore an outcome of
the measured curves, not a hard-coded choice.

% Generated by src/scripts/paper_results.py. Do not edit.
\begin{table*}[t]
\centering\footnotesize
\caption{Where to speculate: the listener against the thinker on this device. Every cell is a scorer's output; sources in RESULTS.md.}
\label{tab:where}
\begin{tabular}{p{0.15\textwidth}p{0.42\textwidth}p{0.36\textwidth}}
\toprule
 & Listener (COMMIT-WL-STOP + gate) & Thinker (SPEC-CONTINUE) \\
\midrule
Survival: how much of the speculation is kept & words committed before the endpoint: 0.409 [0.340, 0.454] n=70 items (STOP + gate, pass 6; P8) & 96 draft tokens per turn; per-item median of accepted / produced 0.000 (n=38 items) (pass 7) \\
Value: TTFA, paired per item, multi\_step & reduction vs REACTIVE-NOPARTIAL +357 ms [+294, +412] n=70 (pass 2k) & difference SPEC-CONTINUE - REACTIVE +32 ms [-92, +92] n=22, positive = slower: no saving detected (pass 7, descriptive) \\
Overlap cost: a decode in flight at the endpoint & final decode +492 ms/s [439, 574] per s of overlap; 960 turns, 43 runs, CI over runs; arms with partial decodes in flight (P2a, not the STOP arm, which freezes them). Partials in flight cost canonical REACTIVE +305.7 ms [+188.0, +452.6] n=15 of TTFA on dev against NOPARTIAL (pass 1n) & TTFA +96.5 ms [+79.4, +109.8] n=258 when the draft is still decoding; +17.5 ms [-24.4, +76.3] n=66 when it finished (blocked policy runs, 2026-09-18/20; the onset test found a finished-draft penalty at one contended cell) \\
Energy per turn, net of the board floor & vs REACTIVE-NOPARTIAL +17.83 J [+17.12, +18.58] n=70 (pass 2k); +0.0499 J/ms [+0.0425, +0.0607] per ms saved; resumed-and-discarded stale decodes mean 1.74 J, total 101.1 J net over 58 decodes & vs REACTIVE +35.3 J [+34.6, +35.7] n=22 (pass 7, descriptive) \\
Memory & the decoding worker process: PSS median 121.4 MB, MemAvailable 118.9 MB lower (pass 2k, at spawn). Keeping an idle second engine resident: +19.1 MB RSS, n=3/3 runs, no detectable TTFA cost +1.8 ms [-30.3, +37.0] n=16 (pass 1r) & no second model: drafts run in the reply LLM's single slot (\texttt{-{}-parallel} 1, fixed context); not measured separately \\
\bottomrule
\end{tabular}
\end{table*}


Table~\ref{tab:where} puts the two consumers side by side on survival, value,
overlap cost, energy and memory.

\section{Is the gate enough? A planner}
\label{sec:planner}

The gate is a single threshold on a point prior. A planner that reasons about
when the user will stop could, in principle, time hypotheses better. We
formalise that planner and evaluate it in simulation. No live pass was run.

\textbf{Formulation.} State is the buffered audio $t$ and the buffer ends
$x_{k-1}$, $x_{k-2}$ of the last two completed hypotheses. Under
LocalAgreement-2 these determine what is committed. At each idle instant the
planner either issues a hypothesis on the buffer or waits. A hypothesis issued
at $t$ completes after its decode time, unless the endpoint arrives first; at
the endpoint it is frozen and lost. The objective is the expected committed
audio at the endpoint. That minimises the expected final decode, because the
final is affine in the uncommitted tail and the freeze removes the overlap
term. The endpoint enters through its hazard, from the smoothed distribution
of utterance lengths. A finite-horizon dynamic program over $t$ gives the
policy table. The planner sees neither the item's length nor its segment
boundaries.

\textbf{Calibration.} The simulator draws, per item, from the logs:
\begin{itemize}
\item the buffer at the endpoint;
\item the segment boundaries the item ever committed, each with its own commit
  rate;
\item partial decode times with their residuals;
\item the final's within-item cost per second of audio.
\end{itemize}
Boundaries that no logged run reached are filled synthetically, at a rate
calibrated so that the simulated committed fractions match those measured with
and without the gate. Run through the simulator, the gate reproduces its
saving as \PlannerVal, against \PthreeRed{} measured; the value measured on
the same items lies just outside the simulated interval. This is a weak
check: the final's cost per second was fitted largely from the same contrast,
and the fill rate was chosen to match the measured committed fractions, so
those fractions are a calibration target rather than a test.

\textbf{Result.} At the gate's duty bound, the median paired per-item TTFA
reduction of the planner over the gate is \PlannerDP. Without the bound it is
\PlannerNoDuty, and that gain comes from decoding back to back, not from
planning. Under the measured costs, the gate is within \PlannerDP{} of the
planner's optimum at the median. The simple rule is near-optimal, which is why
it can be simple. What gain remains sits in items longer than the gate's
prior, where the gate stops issuing early. A gate whose prior adapts to the
elapsed speech might capture it; that is a hypothesis, not simulated here.

\section{Summary}

The listening window goes to the listener: the allocator never spends it on
the thinker, and the listener's gate spends it where commits can pay. With the
gate and a freeze at the endpoint, the listener cuts the median paired TTFA on
long utterances by \PthreeRed{}, abstains on short ones, and costs
\EnergyListenerr{} per turn. On this task the allocator never spends the
window on the thinker, and its log shows why. In simulation, a planner that
reasons about the endpoint's hazard beats the gate by a median of
\PlannerDP{}.
